
\subsubsection{2.1 FAIR Principles and Empirical Validation}

The FAIR Guiding Principles \citep{Wilkinsonetal2016} define four dimensions for scientific data: Findable (persistent identifiers, rich metadata), Accessible (open protocols, clear access conditions), Interoperable (standardized vocabularies, community formats), and Reusable (clear licenses, provenance, usage guidance). Automated FAIR assessment tools such as F-UJI [Devaraju and Huber, 2021] have been developed for general scientific repositories.

Empirical validation of FAIR compliance effects on research outcomes remains sparse. To our knowledge, no systematic large-scale study has established a relationship between FAIR scores and research engagement metrics for ML-specific repositories while controlling for dataset age, prominence, or domain. No prior work applies matched observational designs to isolate FAIR sub-criteria effects on ML repository data, and no study characterizes the specific confounding structure — suppressor confounding by dataset age — that this paper addresses.

\subsubsection{2.2 Dataset Documentation Frameworks}

Datasheets for Datasets \citep{Gebruetal2021} defines documentation dimensions and has been adopted as a reporting standard. Data Cards \citep{Pushkarnaetal2022} operationalize analogous dimensions for HuggingFace. Pineau et al. [2021] linked reproducibility checklists to paper-level outcomes. These frameworks specify what to document but do not provide large-scale empirical evidence that documentation completeness causally increases research adoption. Our work addresses this causal gap for the Findable dimension specifically.

\subsubsection{2.3 Repository Infrastructure and Usage Analysis}

Vanschoren et al. [2014, 2019] introduced OpenML as a collaborative benchmarking platform whose structured run history provides a proxy for deliberate research engagement. Lhoest et al. [2021] describe the HuggingFace Dataset Hub. Prior analyses of OpenML usage patterns have characterized dataset popularity but none have applied FAIR scoring to predict engagement trajectories.

\subsubsection{2.4 Observational Causal Inference}

Propensity score matching [Rosenbaum and Rubin, 1983] creates pseudo-experimental conditions where measured confounders are balanced, enabling causal inference from observational data. Kaplan-Meier survival analysis [Kaplan and Meier, 1958] and Cox proportional hazards regression \citep{Cox1972} provide tools for time-to-event outcomes. Applied to TTFR, these methods capture discovery dynamics rather than static engagement levels, which is appropriate for studying the friction-reduction mechanism that FAIR Findable compliance is hypothesized to provide.

\subsubsection{2.5 Positioning}

Our work differs from prior FAIR-outcome studies by: (1) applying a matched observational design rather than unadjusted correlation, (2) targeting ML-specific repositories with run-count outcomes, (3) disaggregating FAIR sub-criteria, and (4) explicitly characterizing the suppressor confounding structure that is specific to ML repository data.

